% Caso ainda não tenham colocado os pacotes abaizo, eles serão necessários
%\usepackage[utf8]{inputenc}
%\usepackage[T1]{fontenc}
%\usepackage[margin=2.5cm]{geometry}
%\usepackage{enumitem} 
%\usepackage{booktabs} 

%\renewcommand{\baselinestretch}{1.15}

Nesta seção é mostrada a Descrição do Hardware, que são listados os principais componentes que compõem as principais soluções levantadas para o desenvolvimento do projeto. Eles são selecionados com base nas funcionalidades descritas e acordadas com o cliente na etapa de definição de requisitos. 

\subsection{Arquitetura Geral}
O hardware do projeto foi projetado para operar de forma autônoma e centralizada. O sistema baseia-se em um microcontrolador com conectividade Wi-Fi nativa, que atua simultaneamente como gerenciador de aquisição de dados ambientais, controlador de atuadores visuais (LEDs) e hospedeiro do servidor web local. A comunicação entre os módulos periféricos foi distribuída utilizando barramentos padronizados (I2C e SPI) para garantir estabilidade e economia de portas lógicas (GPIOs).

\subsection{Componentes e Justificativas Técnicas}
\begin{itemize}[leftmargin=*]
    \item \textbf{Microcontrolador: ESP32 DevKitC V4 (WROOM-32D)}
    \begin{itemize}
        \item \textit{Especificações:} Processador Dual Core de 32-bits, Wi-Fi 802.11 b/g/n, operação em 3,3V.
        \item \textit{Justificativa:} Escolhido por sua capacidade de processamento multitarefa. O uso de núcleos separados permite manter o servidor web local operando de forma estável, sem interromper as rotinas de leitura de sensores e gravação de logs.
    \end{itemize}

    \item \textbf{Sensor de Temperatura e Umidade: Módulo AHT20}
    \begin{itemize}
        \item \textit{Especificações:} Protocolo I2C, precisão de $\pm 0.3$ $^\circ$C (Temp) e $\pm 2\%$ (Umid), calibração interna de fábrica.
        \item \textit{Justificativa:} Superior em estabilidade aos sensores da família DHT. O uso do barramento I2C minimiza o uso de pinos do microcontrolador e evita erros de temporização comuns em protocolos \textit{single-wire}.
    \end{itemize}

    \item \textbf{Módulo de Armazenamento (Log Local): Leitor Micro SD Padrão}
    \begin{itemize}
        \item \textit{Especificações:} Comunicação SPI, regulador de tensão integrado (5V para 3,3V) e chip \textit{Level Shifter}.
        \item \textit{Justificativa:} O modelo padrão protege o cartão SD e o ESP32 contra picos de tensão. É responsável por gravar as leituras em formato \texttt{.csv/.txt}, garantindo a persistência dos dados em caso de falha de conexão.
    \end{itemize}

    \item \textbf{Interface Visual: LEDs}
    \begin{itemize}
        \item \textit{Justificativa:} Indicam os estados de operação (Ligado, Automático, Desligado). Foram dimensionados em série com resistores de $220\,\Omega$, mantendo o consumo individual em cerca de 6mA e protegendo as portas lógicas do ESP32 contra sobrecorrente.
    \end{itemize}
\end{itemize}

\subsection{Dimensionamento e Alimentação}
O circuito foi projetado para ser alimentado diretamente pela porta USB de um notebook (5V) conectada à entrada Micro-USB do ESP32. O microcontrolador utiliza seu regulador interno (LDO) para converter e distribuir 3,3V de forma segura para os sensores e o cartão SD. O consumo total estimado em pleno funcionamento (Wi-Fi ativo, módulo SD escrevendo e LEDs acesos) não ultrapassa 350mA. Esse valor é perfeitamente suportado pelo limite padrão de 500mA de uma porta USB 2.0 convencional, dispensando o uso de fontes externas durante o desenvolvimento e a apresentação.

\subsection{Estrutura Física e Montagem}
Para atender aos requisitos de robustez do projeto, o circuito final não utilizará protoboard. Todos os componentes serão organizados e soldados em uma placa perfurada ilhada ($5\times7$ cm). O microcontrolador será fixado por meio de barras de pinos fêmea, permitindo sua remoção para manutenção sem a necessidade de dessoldagem.

\subsection{Mapa de Conexões (Pinagem Validada)}
A distribuição das portas foi estruturada consultando o \textit{datasheet} oficial do ESP32 DevKitC V4. Foram rigorosamente evitados pinos classificados como \textit{Input Only} ou \textit{Strapping Pins} (que afetam o boot do sistema) para os atuadores.

\begin{table}[h]
\centering
\caption{Mapeamento Definitivo de Pinagem do Sistema}
\vspace{0.2cm}
\begin{tabular}{@{}llc@{}}
\toprule
\textbf{Componente} & \textbf{Função / Protocolo} & \textbf{Pino do ESP32} \\ \midrule
Sensor AHT20        & I2C (SDA)                   & GPIO 21                \\
Sensor AHT20        & I2C (SCL)                   & GPIO 22                \\ \addlinespace
Módulo SD           & SPI (CS / SS)               & GPIO 5                 \\
Módulo SD           & SPI (MOSI)                  & GPIO 23                \\
Módulo SD           & SPI (MISO)                  & GPIO 19                \\
Módulo SD           & SPI (SCK)                   & GPIO 18                \\ \addlinespace
LED Verde           & Saída Digital (Ligado)      & GPIO 27                \\
LED Amarelo         & Saída Digital (Automático)  & GPIO 26                \\
LED Vermelho        & Saída Digital (Desligado)   & GPIO 25                \\ \bottomrule
\end{tabular}
\end{table}
